% ================================================================
% Scope and Preliminary Design  (answers AC215 MS1 Step 4)
% Needs only packages already in the preamble: tikz (arrows.meta),
% booktabs, graphicx, float, xcolor (dvipsnames), amssymb
% ================================================================
\section*{Scope and Preliminary Design}

TasteBridge serves pairs of readers who each give
$5$ favorite books, using a fixed catalogue of up to $5000$ fantasy and mystery works from the UCSD Goodreads datasets~\cite{wan2018,wan2019}. It is anonymous and has no accounts. Groups of 3--5 are a stretch goal.

\begin{table}[H]
\centering
\small
\begin{tabular}{@{}ll@{}}
\toprule
\textbf{Objective} & \textbf{Success looks like} \\
\midrule
1. Recommend from five books per reader & Beats simple baselines on held-out likes \\
2. Suit both readers, not just one       & Better for the less-satisfied reader in simulated pairs \\
3. Explain each pick to each reader      & Both readers would read $\geq 60\%$ of picks (user study) \\
4. Deploy a reproducible app             & Under 0.5\,s per request; one-command rebuild \\
\bottomrule
\end{tabular}
\caption{Objectives and success criteria.}
\label{tab:objectives}
\end{table}

The project covers three minimum components: large data (tens of millions of
Goodreads interactions), a complex model that learns from ratings and book
content and combines two readers fairly, and scalability through Kubernetes
autoscaling. We chose methods we already understand (implicit-feedback matrix
factorization~\cite{hu2008}, hybrid models for new users~\cite{kula2015}, the
least-misery rule from group recommendation~\cite{dara2020} and
UMAP~\cite{mcinnes2018}), so our new learning can focus on AC215's MLOps tools.
It is also fun to build: we can test it on each other and our friends, and the
galaxy gives us room for a playful front end.
% If the Figma version is attached, add: "A higher-fidelity Figma version of
% screens 3--4 is attached."

